% PS6 mathematical companion. Build with make in this directory.
% The slide order follows the assignment README; no reference results appear here.
% Notation follows lecture L14a, section 'Building a General Factor Model': growth rate g_i(t),
% market g_m(t), economic factors x_k(t), sector factors f_s(t), loadings alpha_i, beta_i, eta_ik,
% delta_is, stock-specific part u_i(t), design matrix \hat{X}, sector matrix B, factor matrix Phi,
% Cholesky factor A, and uncorrelated factors Z.
% The starter docstrings cite five slides by number and title: 6 'One solve for fitting steps 1, 2,
% and 3', 7 'The sector matrix B', 9 'Economic factors from FRED', 10 'Fitting step 4:
% unmix the factors', and 12 'Sampling: draw, mix, and compose'. Update those citations if the slide
% order changes.
% Settings the slides quote. Check each one after an edit of config/defaults.toml or src/Data.jl:
%   steps_per_year = 252 (Δt = 1/252 year): slide 4.
%   375 stocks, 11 sectors, 2,767 price days (T = 2,766): slides 4, 5, and 11.
%   fred_series, fred_kinds, and fred_names (four series): slides 5 and 9.
%   number_of_states = 100, nu = 5.0, epsilon = 1.0e-4, lambda = 100.0: slide 11.
%   number_of_paths = 10 and seed = 5660: slides 12 and 13.
%   benchmark_ticker = "SPY": slides 2, 3, 4, and 5.
\documentclass[aspectratio=169,11pt]{beamer}
\usepackage{vnslides}
\newcommand{\Xh}{\hat{\mathbf{X}}}
\newcommand{\Bb}{\mathbf{B}}
\newcommand{\Fb}{\mathbf{F}}
\newcommand{\Ab}{\mathbf{A}}
\newcommand{\Zb}{\mathbf{Z}}
\newcommand{\Ib}{\mathbf{I}}
\newcommand{\Phib}{\symbf{\Phi}}
\newcommand{\Sigmab}{\symbf{\Sigma}}
\newcommand{\mub}{\symbf{\mu}}
\newcommand{\thetab}{\symbf{\theta}}
\newcommand{\deltab}{\symbf{\delta}}
\newcommand{\epsb}{\symbf{\varepsilon}}
\newcommand{\gb}{\mathbf{g}}
\newcommand{\ub}{\mathbf{u}}
\newcommand{\yb}{\mathbf{y}}
\newcommand{\Xb}{\mathbf{X}}
\newcommand{\fb}{\mathbf{f}}
\newcommand{\onev}{\mathbf{1}}
\newcommand{\syn}{^{\text{sim}}}
% The monospace font has no Greek letters, so code names such as Δt take their glyphs from DejaVu Sans Mono.
\newfontfamily\greekmono{DejaVuSansMono.ttf}[Scale=MatchLowercase]
\newcommand{\gm}[1]{{\greekmono #1}}

\title{\texorpdfstring{{\vnheavy PS6}: Does a Synthetic Market\\Move Together Like\\the Real Market?}{PS6: Does a Synthetic Market Move Together Like the Real Market?}}
\subtitle{CHEME 4/5660 Cornell University | Mathematical Companion}
\author{Copyright Jeffrey Varner 2026. All Rights Reserved}
\begin{document}
\vntitlepage

\begin{frame}{Disclaimer and risks}
\textbf{Training only.} This assignment uses the daily closing prices of SPY
and 375 of its stocks from 2014 to 2024, and four FRED series, for teaching.
This assignment provides no investment advice.

\vspace{1em}
\textbf{Synthetic markets.} A synthetic market is a set of random paths drawn
from a fitted model. A synthetic market is not a forecast of the real market.

\vspace{1em}
\textbf{Simplifying assumptions.} The model uses daily growth rates only. We
omit dividends, commissions, and taxes.
\end{frame}

\begin{frame}{The central question}
Does a synthetic market built from single-asset jump hidden Markov models move
together like the real market? We'll fit the general factor model of lecture
L14a to SPY and 375 of its stocks, draw synthetic markets, and compare how the
stocks move together.
\vspace{0.2em}
\begin{itemize}\setlength{\itemsep}{0.3em}
\item \hilite{Standard: the market and the sectors.} Fitting steps 1, 2, 3,
and 5, with SPY and the 11 sector factors.
\item \hilite{Advanced: economic factors from FRED.} Standard plus four daily
economic factors in fitting step 1.
\item \hilite{Expert: the Cholesky step.} Advanced plus fitting step 4, which
keeps the correlations between the factors in the synthetic market.
\end{itemize}
\slidenote{Reading guide}{All tracks use slides 4–\pageref{fr:residuals} and
\pageref{fr:jump}–\pageref{fr:report}. Slide \pageref{fr:fred} is for
Advanced and Expert, and slide \pageref{fr:unmix} is for Expert only.}
\end{frame}

\begin{frame}{The data and the growth rates}\label{fr:data}
The real market is SPY and 375 of its stocks in 11 GICS sectors, with one
ticker per company. Each stock has a closing price $S_i(t)$ on every trading
day from 2014 to 2024. The daily growth rate of stock $i$, in 1/yr, is:
\[
 g_i(t)=\frac{1}{\Delta t}\ln\left(\frac{S_i(t)}{S_i(t-1)}\right),
 \qquad \Delta t=\tfrac{1}{252}\ \text{yr}.
\]
The 2,767 price days give $T=2{,}766$ growth rates per stock. The growth rates
of SPY are $g_m(t)$, the market. The vector $\gb_i$ holds the $T$ growth rates
of stock $i$.
\slidenote{Code}{\texttt{load\_universe} computes the growth rates with
\texttt{JumpHMM.excess\_growth\_rates(prices; rf = 0.0, \gm{Δt})}, as in the L14a
example.}
\end{frame}

\begin{frame}{The general factor model}\label{fr:model}
The general factor model explains the daily growth rate of stock $i$ with the
market, $K$ economic factors, and the 11 sector factors:
\[
 g_{i}(t)=\underbrace{\alpha_{i}}_{\text{baseline}}
 +\underbrace{\beta_{i}\,g_{m}(t)}_{\text{market}}
 +\underbrace{\sum_{k=1}^{K}\eta_{ik}\,x_{k}(t)}_{\text{economic factors}}
 +\underbrace{\sum_{s=1}^{11}\delta_{is}\,f_{s}(t)}_{\text{sectors}}
 +\underbrace{u_{i}(t)}_{\text{stock-specific}}.
\]
\begin{center}\small
\begin{tabular}{lll}
\toprule
Market in the report & Economic factors & Sector factors \\
\midrule
single index model & none ($K=0$) & none \\
market and sectors (Standard) & none ($K=0$) & 11 \\
market, FRED series, and sectors (Advanced) & four FRED series ($K=4$) & 11 \\
\bottomrule
\end{tabular}
\end{center}
\slidenote{Expert}{The Expert market adds fitting step 4 to the Advanced
market, slide \pageref{fr:unmix}.}
\end{frame}

\begin{frame}{One solve for fitting steps 1, 2, and 3}\label{fr:solve}
For a matrix $\Xb$ with full column rank and a vector $\yb$, the least-squares
solution $\thetab$ minimizes $\lVert\yb-\Xb\thetab\rVert^{2}$. The solution
solves the normal equations:
\[
 \Xb^{\top}\Xb\,\thetab=\Xb^{\top}\yb,
 \qquad
 \thetab=\left(\Xb^{\top}\Xb\right)^{-1}\Xb^{\top}\yb=\Xb^{+}\yb,
\]
where $\Xb^{+}$ is the pseudo-inverse of $\Xb$. The residual $\yb-\Xb\thetab$
is orthogonal to every column of $\Xb$. Fitting steps 1, 2, and 3 use this one
solve with different inputs:
\begin{center}\small
\begin{tabular}{llll}
\toprule
Fitting step & $\Xb$ & $\yb$ & $\thetab$ \\
\midrule
1, for each stock $i$ & $\Xh=[\onev,\gb_m,\mathbf{x}_1,\dots,\mathbf{x}_K]$, $T\times(K+2)$ & $\gb_i$ & $(\alpha_i,\beta_i,\eta_{i1},\dots,\eta_{iK})$ \\
2, for each day $t$ & $\Bb$, $n\times 11$ & $\epsb(t)$ & $\fb(t)=(f_1(t),\dots,f_{11}(t))$ \\
3, for each stock $i$ & $\Fb$, $T\times 11$ & $\epsb_i$ & $\deltab_i=(\delta_{i1},\dots,\delta_{i,11})$ \\
\bottomrule
\end{tabular}
\end{center}
\slidenote{Code}{\texttt{least\_squares\_fit(X, y)} returns $\thetab$ for each
of the three steps.}
\end{frame}

\begin{frame}{The sector matrix B}\label{fr:sector}
The sector matrix $\Bb$ has one row for each of the $n=375$ stocks and one
column for each sector:
\[
 B_{is}=\begin{cases}1 & \text{if stock } i \text{ is in sector } s,\\
 0 & \text{otherwise.}\end{cases}
\]
Each row of $\Bb$ holds exactly one 1. The matrix $\Bb^{\top}\Bb$ is diagonal,
with the number of stocks $n_s$ of each sector on the diagonal, and entry $s$
of $\Bb^{\top}\epsb(t)$ is the sum of the residuals of the stocks in sector $s$
on day $t$. Fitting step 2 therefore gives the average residual of each sector:
\[
 f_s(t)=\frac{1}{n_s}\sum_{i\in s}\varepsilon_i(t).
\]
\slidenote{Code}{\texttt{sector\_matrix(sectors, sector\_names)} returns
$\Bb$, with the columns in the order of \texttt{sector\_names}.}
\end{frame}

\begin{frame}{Fitting steps 1, 2, and 3: the residuals}\label{fr:residuals}
Fitting step 1 leaves the residual $\epsb_i=\gb_i-\Xh\thetab_i$ of each stock.
Fitting step 2 stacks the sector factors $\fb(t)$ as the rows of the
$T\times 11$ matrix $\Fb$. Fitting step 3 leaves the stock-specific part:
\[
 \ub_i=\epsb_i-\Fb\deltab_i,
 \qquad\text{so}\qquad
 \gb_i=\Xh\thetab_i+\Fb\deltab_i+\ub_i
\]
holds exactly on every day. Each residual is orthogonal to the columns of its
own $\Xb$, so each $\ub_i$ is uncorrelated with the market, the economic
factors, and the sector factors.
\slidenote{Single index model}{The single index model stops after fitting step
1 with $K=0$. Its stock-specific part is $\ub_i=\epsb_i$, and its $\Fb$ has
zero columns.}
\end{frame}

\begin{frame}{Economic factors from FRED}\label{fr:fred}
FRED gives the level $L(t)$ of each series on each price day. A rate becomes
its daily change, and a price becomes its growth rate, in 1/yr:
\[
 x(t)=L(t)-L(t-1)\ \ \text{(change)},
 \qquad
 x(t)=\frac{1}{\Delta t}\ln\left(\frac{L(t)}{L(t-1)}\right)\ \ \text{(growth)}.
\]
\begin{center}\small
\begin{tabular}{lll}
\toprule
FRED series & Level & Kind \\
\midrule
\texttt{DGS10} & 10-year Treasury yield (percent) & change \\
\texttt{T10YIE} & 10-year breakeven inflation rate (percent) & change \\
\texttt{DCOILBRENTEU} & Brent crude oil price (USD per barrel) & growth \\
\texttt{DTWEXBGS} & US dollar index against a broad basket of currencies & growth \\
\bottomrule
\end{tabular}
\end{center}
\slidenote{Code}{\texttt{economic\_factor(levels, kind, \gm{Δt})} returns
$x(1),\dots,x(T)$. This product uses the FRED\textsuperscript{\textregistered}
API but is not endorsed or certified by the Federal Reserve Bank of St.\ Louis.}
\end{frame}

\begin{frame}{Fitting step 4: unmix the factors}\label{fr:unmix}
The factor matrix $\Phib=[\gb_m,\mathbf{x}_1,\dots,\mathbf{x}_K,\Fb]$ is
$T\times p$, with $p=1+K+11$. Its column averages are $\mub$, and its sample
covariance matrix is:
\[
 \Sigmab_{\Phi}=\frac{1}{T-1}\left(\Phib-\onev\mub^{\top}\right)^{\top}\left(\Phib-\onev\mub^{\top}\right).
\]
The Cholesky factorization $\Sigmab_{\Phi}=\Ab\Ab^{\top}$ gives the
lower-triangular matrix $\Ab$. The uncorrelated factors are:
\[
 \Zb=\left(\Phib-\onev\mub^{\top}\right)\Ab^{-\top},
 \qquad
 \operatorname{cov}(\Zb)=\Ab^{-1}\Sigmab_{\Phi}\Ab^{-\top}=\Ab^{-1}\Ab\Ab^{\top}\Ab^{-\top}=\Ib.
\]
\slidenote{Code}{\texttt{unmix\_factors(\gm{Φ})} returns $(\mub,\Ab,\Zb)$. In
Julia, \texttt{M / A'} computes $\mathbf{M}\Ab^{-\top}$.}
\end{frame}

\begin{frame}{Fitting step 5: the jump hidden Markov models}\label{fr:jump}
Fitting step 5 fits one single-asset jump hidden Markov model to each series:
\begin{itemize}\setlength{\itemsep}{0.3em}
\item each column of $\Phib$ without fitting step 4 (Standard and Advanced), or
each column of $\Zb$ with fitting step 4 (Expert);
\item the stock-specific part $\ub_i$ of each of the 375 stocks.
\end{itemize}
Every model has 100 hidden states, a Student-$t$ distribution with $\nu=5$ in
each state, and the jump settings $\epsilon=10^{-4}$ and $\lambda=100$ days of
the tuned SPY model of L13a. Each model draws its paths independently of the
other models. Without fitting step 4, the synthetic factors are uncorrelated,
so the covariance of two synthetic stocks keeps only the variance of each
factor, and every covariance term between two different factors drops out.
\slidenote{Code}{\texttt{fit\_jump\_model} in \texttt{src/Model.jl} calls
\texttt{JumpHMM.build}; \texttt{draw\_series} calls \texttt{JumpHMM.simulate}.
FRED records the 10-year yield and breakeven inflation to 0.01 percentage
point, so \texttt{spread\_rate\_changes} spreads each daily change over its
0.01 step before the model of that factor is fit.}
\end{frame}

\begin{frame}{Sampling: draw, mix, and compose}\label{fr:sampling}
Each synthetic market has $T$ days, as long as the data:
\begin{enumerate}\setlength{\itemsep}{0.3em}
\item \hilite{Factor draws.} Draw a path of $T$ days from each factor model.
With fitting step 4, the draws are the columns of $\Zb\syn$, and we mix the
draws back:
\[
 \Phib\syn=\onev\mub^{\top}+\Zb\syn\Ab^{\top},
 \qquad
 \operatorname{cov}(\Phib\syn)=\Ab\,\operatorname{cov}(\Zb\syn)\,\Ab^{\top}\approx\Ab\Ab^{\top}=\Sigmab_{\Phi}.
\]
Without fitting step 4, the draws are the columns of $\Phib\syn$.
\item \hilite{Stock-specific draws.} Draw a path $\ub_i\syn$ from the model of
each $\ub_i$.
\item \hilite{Compose.} With $\Xh\syn$ and $\Fb\syn$ taken from the columns of
$\Phib\syn$: $\gb_i\syn=\Xh\syn\thetab_i+\Fb\syn\deltab_i+\ub_i\syn$.
\end{enumerate}
\slidenote{Code}{\texttt{mix\_factors(Z, \gm{μ}, A)} returns $\Phib\syn$, and
\texttt{compose\_growth\_rates(X, \gm{θ}, F, \gm{δ}, u)} returns $\gb_i\syn$. The
checker draws 10 synthetic markets from the seed 5660.}
\end{frame}

\begin{frame}{The columns of the report}\label{fr:report}
With $\rho_{ij}$ the correlation of the growth rates of stocks $i$ and $j$ over
the $T$ days, and $s(\cdot)$ the sample standard deviation:
\vspace{0.3em}
\begin{center}\small
\begin{tabular}{ll}
\toprule
Column & Definition \\
\midrule
Same sector & average $\rho_{ij}$ over pairs $i<j$ in the same sector \\
Different sectors & average $\rho_{ij}$ over pairs $i<j$ in different sectors \\
Each pair of sectors & average $\rho_{ij}$ over $i$ in the first sector and $j$ in the second \\
Median stock volatility (1/$\sqrt{\text{yr}}$) & median over the stocks of $s(\gb_i)\sqrt{\Delta t}$ \\
R² (table FRED series) & $1-\operatorname{var}(\epsb_i)/\operatorname{var}(\gb_i)$, after fitting step 1 \\
\bottomrule
\end{tabular}
\end{center}
\vspace{0.3em}
Each synthetic row averages the 10 synthetic markets. The table FRED series
averages R² and each loading $\eta_{ik}$ over the stocks of each sector.
\end{frame}

\begin{frame}{Summary}
PS6 fits the general factor model of lecture L14a to SPY and 375 of its
stocks, draws synthetic markets from single-asset jump hidden Markov models,
and compares how the stocks move together.

\vspace{0.6em}
\textbf{Key takeaways}
\begin{itemize}\setlength{\itemsep}{0.5em}
\item Fitting steps 1, 2, and 3 are one least-squares solve,
$\Xb^{\top}\Xb\,\thetab=\Xb^{\top}\yb$, with a different $\Xb$ and $\yb$ in each
step. With $\Xb=\Bb$, the solve gives the average residual of each sector.
\item Separate jump hidden Markov models draw independent paths, so the
synthetic factors are uncorrelated unless the factors are unmixed first.
\item The Cholesky factor $\Ab$ of the covariance matrix unmixes the factors
before the fit and mixes the draws back after sampling.
\end{itemize}
\vspace{0.6em}
Use these steps to answer the central question.
\end{frame}
\end{document}
